\chapter{Conclusions and Future Work}
\label{sec:conclusion}

\section{Summary of findings}
\label{sec:concl-summary}

This thesis asked whether a compact language model, specialised in a single decision and used as a decision model, can
verify the answers of a RAG system sentence by sentence on a laptop, and where any advantage comes from. GroundOne, a
fine-tuned Qwen2.5 model that reads the verdicts of a whole answer in a single forward pass, gives a positive answer,
albeit with some reservations.

Specialisation works where it was aimed. On the RAGTruth test set GroundOne reaches a three-way balanced accuracy of
0.769 with 0.5B parameters and 0.832 with 3B. It is far above the same models untrained, ranks supported sentences better
than the specialised encoders and is clearly above, in three-way balanced accuracy, Rizzo-Flow, a general-purpose
decision model eight times larger than the 0.5B. On LLM-AggreFact the 3B is on par with the best small encoder. On
VitaminC, instead, the gain disappears: a specialist is not a general fact-checking tool. d1-3B, a general-purpose
decision model of the same size trained to decide on a broad mixture of tasks, confirms the trade-off: GroundOne is far
ahead of it on RAGTruth, with the 3B it matches d1-3B on LLM-AggreFact, and it is far behind on VitaminC.

The decision-model mechanism pays off on two fronts. The first is speed. Reading all the verdicts in a single forward
pass makes the verification of long answers about two times faster than that of a terse LLM judge and six times
faster than that of a verbose LLM judge of the same size, with output that is always valid. The second is accuracy:
restricting the decision to the admissible answers is, by itself, more accurate than letting the same model generate the
verdict. Moreover, even when asked for the same output, LLM judges cannot provide probabilities of comparable
quality. Their probabilities barely separate correct from incorrect decisions, and the verbalised ones are almost always
all-or-nothing.

The probabilities of GroundOne can be calibrated with a temperature and a bias fitted on the dev set. A threshold chosen
on the same data then turns them into labels consistent both with the probabilities and with the evaluated decision.
The Expected Calibration Error alone, however, does not distinguish useful probabilities from useless ones:
discrimination must be measured separately.

The whole system runs on a 16\,GB laptop. A long answer is verified in 0.41 seconds with the 0.5B and in 2.57 with the
3B, with an energy of 5.5 and 33 joules on average per answer; when the document is already in the KV cache, as in the verification of a rewritten answer, 0.11 and 0.58 seconds suffice. Compared with an LLM judge the energy saving is
real, but smaller than the time saving, because a single long prefill draws more power than decoding.

The GroundOne Studio application shows that a complete RAG assistant can be built around the verifier, in a single
process and without network access: retrieval, generation, verification and comparison with an LLM judge require
at most 6.2\,GB of memory allocated by MLX. In the injected-error case study GroundOne flags the error in 97 of the 98
answers in which the sentence receives a verdict, against 84 for the LLM judge, and targeted rewriting reduces the answers
that still contain the error from 100\,\% to 9\,\%, according to the automatic count: numbers and dates are mostly
restored from the source, fabricated sentences deleted. The case study also shows that the type of error matters: a
contradicted sentence is corrected, a not-mentioned one often merely removed. The answers that pass the loop are
accepted by GroundOne, not independently verified.

The main limitations are two. The first is the narrowness of the training data, because of which contradictions outside
data-to-text and hallucinations hidden in copied sentences are learned only in part. The second is a training recipe that
is not fully robust to the choice of the seed.

\section{Future work}
\label{sec:concl-future}

Several lines of work follow directly from these limitations. The training recipe should be made robust by enabling bias
correction and gradient clipping, logging the gradient norms and repeating the runs over several seeds, including for the
3B. Training data with more varied contradictions, with paraphrased support and with minimal edits of the kind VitaminC
contains would then be needed: they should address both the low recall outside data-to-text and the blind spot on copied
sentences. Finally, training the readout on the letter that the model would spontaneously write after the marker could
align readout and generation, with effects on accuracy that remain to be measured.

On the systems side, the cost model of Chapter~\ref{sec:local} should be verified on other hardware. A port of the
readout to llama.cpp would allow GroundOne to be measured on a discrete graphics card, on a DGX Spark or an AMD Strix Halo
machine, and on a smartphone, whose neural processing units accelerate precisely the prefill on which a decision model
relies. Latency comparisons with the other families of verifiers also require fair conditions: the encoders must be run
in batches, while GroundOne and the general-purpose decision models must be queried through the same HTTP interface. Two changes would also shorten the reading of the text, which is now almost all of GroundOne's time. The first is fusing the LoRA adapter into the base weights, which would remove the roughly 10\,\% of extra computation measured in the prefill, at the price of a new quantisation and a new calibration. The second is streaming verification: reading the document while the generator writes and appending the sentences as they arrive, using prefix reuse, so that at the end of generation only the last sentence remains to be read.

Another direction concerns language. The first step is to measure how well the adapter trained in English already works
on documents and answers in another language, for example by translating part of T1 into Italian. If the transfer is not
sufficient, a model this small may be better served by a solution other than a single multilingual adapter: a twin
adapter for each language, trained separately on the same base, so that the languages do not have to share the little
capacity available. Adapters take up little space and can be loaded and unloaded at runtime on top of a single copy of
the base model, with a language identifier choosing the right one. Three problems remain to be solved: labelled data are
needed in every language, because translating RAGTruth breaks the alignment of the spans and requires the translated
labels to be checked; temperature, bias and threshold must be fitted again for each language; and the frequent case of a
document in one language and an answer in another falls within no single adapter. The comparison between per-language
adapters and a multilingual adapter is therefore an experiment to be run, not a conclusion.

For contexts made of several documents, a natural extension is attribution: reading, besides the verdict, which passage
supports or contradicts each sentence, for example with a readout per passage, so as to show the user the source of every
claim and to handle conflicting sources explicitly.

The application of Chapter~\ref{ch:studio} contains a first prototype of this attribution, which remains to be evaluated
with data indicating which passage supports each sentence. Three further directions follow from the application. A larger
generator, for example a 7B at 4 bits, would still fit in the memory of the laptop and should write more complete answers
and correct the flagged sentences better. The rewritten sentences should be checked independently, with a second verifier
or with a human reading of a sample, because their correctness is currently judged by GroundOne alone. Verification,
finally, could proceed while the answer is being generated, flagging each sentence as soon as it is complete: the
marker-based design makes this use natural.
